\chapter{INTRODUCTION}\label{ch:introduction}

\section{Overview}\label{sec:overview}
Modern warfare has transitioned into an increasingly data-driven environment where split-second
decisions are critical for survival and mission success. The sheer volume of visual information from
surveillance feeds, drones, and tactical environments often overwhelms personnel, leading to
cognitive fatigue and a higher margin for error. Manual threat assessment is no longer sufficient to
ensure the safety and tactical superiority of ground units in high-stress scenarios. Consequently,
there is an urgent need for intelligent, autonomous systems that can process vast amounts of
real-time visual data to provide actionable insights and enhance situational awareness. \vspace{6pt}

VeerDrishti is an AI-enabled combat vision system designed to bridge this gap by providing soldiers
with a persistent, intelligent observer through a camera-integrated solution. By leveraging
state-of-the-art computer vision models and edge computing, the system autonomously identifies enemy
combatants, vehicles, and weapons. The primary objective is to offload the mental burden of constant
monitoring from the soldier, allowing them to focus on strategic execution. Developed as a
lightweight solution suitable for deployment on wearable hardware, the system aims to transform how
visual data is utilized on the front lines of national security. \vspace{6pt}

The project incorporates sophisticated methodologies to ensure reliability under diverse and
challenging battlefield conditions. It utilizes multimodal image fusion—combining visible and
thermal spectrums—to maintain high detection accuracy in low-light or obscured environments.
Furthermore, by implementing edge-optimized architectures and advanced feature fusion techniques,
the system achieves high-speed inference without compromising precision. This integration of
versatile object detection and robust information processing ensures that the system remains a
reliable asset, offering a significant technological advantage in modern combat and defense
operations.

\label{end:overview}
\section{Problem Statement}\label{sec:problem}
The problem statement is to address the challenges of managing vast real-time visual data in modern
battlefield environments, where manual monitoring is slow, error-prone, and inefficient; in order to
enhance situational awareness and tactical decision-making, an AI-powered combat vision system is
proposed to automatically detect, classify, and assess threats in real time, improving response
speed and operational efficiency.

\label{end:problem}
\section{Motivation}\label{sec:motivation}
The motivation for VeerDrishti is to minimize casualties and save lives by providing an intelligent,
persistent observer in combat environments. Modern warfare involves overwhelming visual data that
leads to cognitive fatigue, increasing the risk of fatal human error. The system serves as a
technological safeguard, enhancing personnel safety and protecting national security forces during
high-stress operations.

\subsection{Scope}
The scope of the project is to develop an autonomous, real-time vision system designed to detect
enemy personnel, weapons, and vehicles in tactical environments. The technical implementation will
focus on developing a lightweight system, optimized for localized inference on mobile edge devices.

\subsection{Purpose}
The purpose of the project is to provide an autonomous, intelligent observer that delivers real-time
situational awareness in high-stress combat environments. By automating threat detection and
mitigating cognitive fatigue, the system serves as a technological safeguard dedicated to achieving
the project's core mission: minimizing casualties and saving lives.

\subsection{Applicability}
The project will be applicable in active combat operations, border surveillance, and high-risk
tactical missions where ground units require instantaneous threat identification. The edge-optimized
architecture will ensure reliability in signal-denied zones and low-visibility environments, serving
as a critical force multiplier for national security personnel. \vspace{6pt}

Beyond defense, the system's ability to automate hazard detection will be directly transferable to
emergency response and disaster management, fulfilling the mandate to minimize casualties and save
lives.

\label{end:motivation}
\section{Objectives}\label{sec:objectives}
The project's primary objectives are to develop a lightweight, real-time combat vision system that
can be deployed on wearable hardware, capable of autonomously detecting and classifying threats in
dynamic battlefield environments. \vspace{6pt}

Specific objectives for the project include:
\begin{enumerate}[label=\roman*.]
  \item Implement a lightweight real-time object detection model.
  \item Enable on-screen visual highlighting of identified threats.
  \item Optimize system performance for low latency and smooth operation.
  \item Evaluate accuracy and response time under simulated field conditions.
  \label{end:objectives}
\end{enumerate}
\section{Methodology}\label{sec:methodology}
The project adopts a hybrid deep-learning approach that integrates generative data augmentation with
edge-optimized multimodal fusion. Initially, the system utilizes a Generative Adversarial Network
(GAN) to synthesize high-quality military training data, addressing scarcity in combat-specific
datasets. For real-time execution, the project implements a lightweight detection architecture, such
as EfficientDet or YOLOv10, optimized for low-latency performance on edge devices like smartphones
or Raspberry Pi. To ensure reliability in varying field conditions, the methodology incorporates
multimodal image fusion, specifically merging RGB texture with Infrared (thermal) edge data using a
fine-grained region attention mechanism. This ensures robust target classification even during night
operations or in cluttered environments. Furthermore, the system leverages a vision-language model
(like Florence v2) for zero-shot detection, allowing the identification of "unseen" or novel threats
through natural language supervision. Post-detection, the pipeline filters data and applies an
Augmented Reality (AR) overlay to highlight threats instantly, while securing all mission metadata
in a local storage layer for subsequent auditing.

\label{end:methodology}
\section{Sustainable Development Goal Relevance}\label{sec:sdg}
The project aligns with United Nations Sustainable Development Goal 16: Peace, Justice, and Strong
Institutions.

\textbf{Target 16.1: Significantly reduce all forms of violence and related death rates everywhere}

VeerDrishti will enhance battlefield situational awareness by detecting and tracking potential
threats in real time. This capability will enable defense personnel to respond faster and make more
informed decisions during combat situations. By improving awareness and reaction time, the system
will help reduce operational risks and potentially decrease casualties. \\

\textbf{Target 16.3: Promote the rule of law at the national and international levels and ensure
  equal access to justice for all}

VeerDrishti will record visual data during operations, enabling improved mission analysis and
operational review. It will support evidence-based evaluation of decisions made during missions.
Such recorded data could also assist investigations related to violations of international
humanitarian law, including potential war crimes. \\

\textbf{Target 16.A: Strengthen relevant national institutions, including through international
  cooperation, for building capacity at all levels, in particular in developing countries, to
  prevent violence and combat terrorism and crime}

VeerDrishti will integrate AI-based detection and tracking technologies, strengthening the
technological capacity of national security institutions. The project will support peacekeeping
operations and the protection of national interests.

\label{end:sdg}
\section{Feasibility Study}\label{sec:feasibility}
The feasibility of the proposed combat vision system is assessed across technical, operational, and
economic dimensions.

\subsection{Technical Feasibility}
The proposed system is technically feasible using existing smartphone hardware and computer vision
frameworks. Modern phones provide high-resolution cameras, sufficient processing power, and GPU
support for real-time object detection. Lightweight AI models can be optimized for edge deployment,
enabling on-device threat detection without requiring complex external hardware.

\subsection{Operational Feasibility}
The system can be integrated into existing helmet setups with minimal modification, using a mounted
smartphone as both camera and display interface. It reduces cognitive load by automating threat
detection and visual highlighting. Basic training would enable personnel to interpret AR overlays
effectively in dynamic environments.

\subsection{Economic Feasibility}
Using commercially available smartphones significantly lowers development and deployment costs
compared to specialized military hardware. Open-source AI frameworks further reduce software
expenses. The prototype approach minimizes initial investment, making the system cost-effective for
research, testing, and potential scalable improvements in future versions.

\label{end:feasibility}
\section{Work Distribution}\label{sec:work}
\begin{table}[h]
  \centering
  \caption{Work Distribution Among Team Members}\label{tab:work}
  \begin{tabular}{lll}
    \toprule
    \textbf{Name} & \textbf{Tasks} \\
    \midrule
    Gaurav Chauhan & Literature Survey and Report Writing \\
    Gaurav Singh & Literature Survey and Slide Preparation \\
    \bottomrule
  \end{tabular}
\end{table}
\label{end:work}\label{end:introduction}
